\documentclass{../latex/utmreport}
\newcommand{\ReportTitle}{Lab 4: The xv6 Shell from the Outside}
% Shared title-page and PDF metadata for every laboratory report.
\newcommand{\StudentName}{Gafenco Victor}
\newcommand{\StudentGroup}{FAF-241}
\newcommand{\CourseName}{Operating Systems}
\newcommand{\InstructorName}{Patricia Reițman}
\newcommand{\InstructorTitle}{Assistant lecturer}
\newcommand{\ReportCity}{Chișinău}
\newcommand{\ReportYear}{2026}
\newcommand{\CourseRepository}{https://github.com/yoda-dinmd/os-labs}

\selectlanguage{english}
\utmlogo{../latex/assets/UTM_Logo_English.jpeg}
\institution{Technical University of Moldova}
\faculty{Faculty of Computers, Informatics, and Microelectronics}
\department{Department of Software Engineering}
\studyprogramme{Software Engineering Study Programme}
\documenttype{Laboratory report}
\student{\StudentName}{\StudentGroup}
\supervisor{\InstructorName}{\InstructorTitle}
\renewcommand{\UTMCoordinatorLabel}{Instructor}
\cityyear{\ReportCity}{\ReportYear}
\reportsubtitle{\CourseName}
\hypersetup{pdfauthor={\StudentName},pdfsubject={\CourseName},
  pdftitle={\ReportTitle}}
\reporttitle{\ReportTitle}
\addto\captionsenglish{\renewcommand{\figurename}{Figure}\renewcommand{\tablename}{Table}\renewcommand{\bibname}{References}}
\renewcommand{\figurename}{Figure}
\renewcommand{\tablename}{Table}
\renewcommand{\bibname}{References}
\usepackage{float}
\usepackage{placeins}
\setlength{\emergencystretch}{2em}
\newcommand{\code}[1]{\texttt{\detokenize{#1}}}
% Keep screenshots beside their observations, with bounded height and width.
\newcommand{\screenshot}[4][\textwidth]{%
  \begin{figure}[H]
    \centering
    \includegraphics[width=#1,height=0.62\textheight,keepaspectratio]{screenshots/#2}
    \caption{#3}\label{#4}
  \end{figure}}
\newcommand{\reportreferences}[2]{%
  \clearpage\phantomsection\addcontentsline{toc}{chapter}{References}
  \begin{thebibliography}{9}
    \bibitem{handout} Technical University of Moldova, FCIM / FAF.
    \textit{#1}. Operating Systems laboratory handout, academic year 2026--2027.
    Supplied course document: \code{#2}.
  \end{thebibliography}}

\begin{document}
\makeutmreporttitlepage
\tableofcontents
\clearpage
\begin{utmabstract}[Abstract]
This laboratory compares a production command interpreter with the small shell
of xv6. A pinned xv6 revision was built on openSUSE Tumbleweed using a RISC-V
cross-compiler and booted in QEMU. A sleep utility and an argument-inspection
program were registered and tested. Parallel Bash and xv6 experiments establish
which facilities are shared and which require additional shell interpretation.
Reading the source explains the absence of expansion, the treatment of quotes,
and a limitation of directory changes inside compound command lines.
\utmkeywords{xv6, Bash, RISC-V, exec, argument vector, shell parsing}
\end{utmabstract}
\utmintroduction
The objective is to distinguish kernel mechanisms from shell conveniences by
observing the same command text under two interpreters. The handout assigns
four parts: sleep compatibility, a shell comparison, inspection of argument
vectors, and source analysis~\cite{handout}.

The host is openSUSE Tumbleweed. Revision
\code{06aad25c735fd3159bdfae5680be4eab7b1668b2} of xv6 is checked out on branch
\code{lab} in \code{~/xv6-riscv}. GCC 14.4.0 with the
\code{riscv64-suse-linux-} prefix and QEMU 11.1.2 provide the build and execution
environment. The course preparation check built the kernel and filesystem and
verified a successful boot and directory listing. Figures show saved terminal
output; the corresponding plain text is retained with the laboratory files.

\chapter{Adapting and testing the sleep utility}
\section{Interface and build integration}
The user header declares \code{int pause(int);} at line 25. The source retained
from Lab 2 already calls this function, so the correct adaptation is to preserve
that call while keeping the executable's name \code{sleep}. The program validates
that exactly one operand is supplied, converts it with \code{atoi}, calls
\code{pause}, and exits. Incorrect argument counts produce a diagnostic on
standard error and terminate with status 1.
\screenshot{1.1-api.png}{The pause declaration used by this xv6 revision}{fig:l4-api}
The Makefile registers both \code{$U/_sleep} and \code{$U/_args} as tab-indented
UPROGS entries with continuation backslashes. A blank line still terminates the
continued list. An additional compiler-detection branch recognizes openSUSE's
prefix, allowing the ordinary \code{make qemu} command to work without a manually
exported variable. The complete Makefile and a focused patch accompany the report.
\screenshot{1.2-sleep-source.png}{Sleep source and its UPROGS registration}{fig:l4-sleep-source}

\section{Execution results and limits}
The command \code{sleep 20} returned normally to the xv6 prompt. Both
\code{sleep} and \code{sleep 10 20} printed \code{usage: sleep <ticks>}.
These observations establish executable availability and argument-count handling.
They do not independently measure the exact waiting duration or display the
process exit status. The implementation follows the exercise's numeric-conversion
scope and does not fully validate malformed tick strings.
\screenshot{1.3-sleep.png}{A valid call and two invalid argument counts}{fig:l4-sleep-test}

\chapter{Comparing the two shells}
\section{Method and results}
Bash ran in \code{~/shtest}, with README and the Linux argument-program source
available. xv6 ran from the root of its filesystem image, where the executables
are installed. Both interpreters received the prescribed command lines. Differences
in listing format are separated from differences in shell interpretation.
\begingroup\small\setstretch{1}
\begin{longtable}{p{0.31\textwidth}p{0.27\textwidth}p{0.33\textwidth}}
\textbf{Input} & \textbf{Bash result} & \textbf{xv6 result}\\\hline
\code{echo hello > f ; cat < f} & hello & hello\\
\code{ls | grep READ} & README & README with type, inode and size columns\\
\code{(echo a; echo b) | wc} & 2 lines, 2 words, 4 bytes & Same counts\\
\code{echo one ; echo two} & one and two on separate lines & Same\\
\code{echo bg &} & Job/PID notice and bg & bg; no Bash-style job notice\\
\code{ls *} & Expanded directory entries & \code{ls: cannot open *}\\
\code{echo $HOME} & /home/victor & Literal \code{$HOME}\\
\code{echo "a b"} & a b without quotes & \code{"a b"} including quotes\\
\code{mkdir d ; cd d ; ls ..} & Changes into d, lists parent & \code{exec cd failed}, then root listing\\
Up arrow, Tab, Ctrl+R & History, completion, reverse search & Raw input bytes; no equivalent editing action\\
\end{longtable}
\endgroup
\screenshot{2.1-bash.png}{The comparison command sequence under Bash}{fig:l4-bash}
\screenshot{2.2-xv6.png}{The same command sequence under xv6}{fig:l4-xv6}

\section{Directory changes and executable search}
The compound \code{mkdir d ; cd d ; ls ..} test reveals an important qualification
to the handout. In this xv6 shell, \code{cd} is recognized only when it begins the
whole input line before normal parsing. Embedded in a semicolon list, it is treated
as an ordinary external executable and fails. Consequently the exact combined
command never changes directory in xv6.

A subsequent standalone \code{cd d} succeeds. From that directory,
\code{ls ..} fails with \code{exec ls failed}, while \code{/ls ..} works.
There is no PATH search: the executable name \code{ls} is relative to the current
directory, whereas \code{/ls} identifies the installed program in the root.
The relative pathname \code{../ls} would also locate it from \code{/d}.
The test ends with a standalone \code{cd /}.
\screenshot{2.3-path.png}{Standalone directory change and the absolute-path workaround}{fig:l4-path}

\section{Shared mechanisms and missing conveniences}
Both shells provide process execution, input/output redirection, pipes, sequencing,
grouping and background execution. Each can change its own current directory,
although xv6's recognition is restricted. Bash additionally provides variable and
pathname expansion, meaningful quote syntax, executable search, history, completion,
conditional operators and job-control commands.

History and completion would be especially noticeable omissions in daily use.
A fixed executable-directory fallback would be a relatively small first addition:
after a suitable failed relative execution attempt, try a known program directory
while retaining the semantics of explicit pathnames. Quote-aware parsing and general
globbing would involve more extensive interpretation. This is an implementation
assessment rather than a modification made in this laboratory.

\chapter{Inspecting the argument vector}
\section{Two implementations of the same observation}
The xv6 program includes its user headers, prints \code{argc}, prints each
\code{argv[i]} inside brackets and exits through \code{exit(0)}. The Linux version
uses \code{stdio.h} and returns zero. Neither version performs globbing, variable
expansion or quote processing. The Linux build succeeds with
\code{gcc -Wall -Wextra -Werror -o args args_linux.c}.
\screenshot{3.1-args-source.png}{Argument-inspection source and build registration}{fig:l4-args-source}

\section{The wildcard and quoted-space experiment}
In the Bash test directory, the matching C filename is \code{args_linux.c}.
The two executions produce the following complete argument vectors:
\begin{lstlisting}
LINUX$ ./args *.c "x y"
argc = 3
argv[0] = [./args]
argv[1] = [args_linux.c]
argv[2] = [x y]

$ args *.c "x y"
argc = 4
argv[0] = [args]
argv[1] = [*.c]
argv[2] = ["x]
argv[3] = [y"]
\end{lstlisting}
Bash passes two operands plus the program name, whereas xv6 passes three operands
plus the program name. The number of Bash operands would increase if more C files
matched. The quoted space survives inside one Bash argument, but xv6 separates at
that space and retains the quote characters in the adjacent words.

\section{Other argument-boundary tests}
\begingroup\small\setstretch{1}
\noindent\begin{tabularx}{\textwidth}{p{0.29\textwidth}YY}
\textbf{Operands typed} & \textbf{Bash} & \textbf{xv6}\\\hline
\code{one    two    three} & argc 4; three words & argc 4; same words\\
\code{$HOME} & argc 2; expanded home pathname & argc 2; literal dollar-name\\
\code{'single quoted'} & argc 2; one unquoted string & argc 3; two strings retaining quotes\\
\end{tabularx}\par
\endgroup
\screenshot{3.2-xv6-args.png}{All argument-vector tests in xv6}{fig:l4-args-xv6}
\screenshot{3.3-linux-args.png}{All argument-vector tests in Bash}{fig:l4-args-bash}
Repeated unquoted whitespace separates words in both interpreters and does not
create empty arguments. The shell constructs the vector before the executable
begins. The kernel's execution interface transfers that vector to the new program;
it does not implement Bash's wildcard or quoting language. The receiving program
merely exposes the values that it was given.

\chapter{Explaining the results from the source}
\section{The five command kinds}
The pinned \code{user/sh.c} contains 499 lines. Lines 8--12 define five node kinds:
\begin{description}[leftmargin=1.25cm]
\item[EXEC] An ordinary command word produces an execution node; no special
operator is required. The execution case invokes \code{exec} with its argument vector.
\item[REDIR] The operators \code{<}, \code{>} and the parsed \code{>>} token produce
redirection nodes. A descriptor is closed and reopened before the nested command runs.
\item[PIPE] The \code{|} operator produces a pipe node. Children connect their input
and output using the pipe descriptors and \code{dup}.
\item[LIST] A semicolon produces a sequence node. The left side is forked and waited
for before the right side runs.
\item[BACK] An ampersand produces a background node, which forks its nested command
without waiting for that child in the node itself.
\end{description}
Parentheses invoke \code{parseblock} to group existing node kinds; they are not a
sixth kind. Recognizing a token does not imply full Bash semantics: xv6's \code{>>}
case opens with \code{O_WRONLY | O_CREATE} without an append flag or a seek, so it
must not be assumed to append like Bash.
\screenshot{4.1-defines.png}{Revision, source length and command-node definitions}{fig:l4-defines}
\screenshot{4.2-cases.png}{Execution and parsing switch cases}{fig:l4-cases}

\section{Why cd belongs to the parent shell}
Lines 166--170 handle \code{cd} directly in \code{main}, before the ordinary fork
and parse path. The persistent shell process must call \code{chdir}; a separate
child's directory change cannot update its parent. The position of this check
also accounts for the failure of an embedded \code{cd} in the earlier list.
\screenshot{4.3-cd.png}{Directory changes handled before forking and parsing}{fig:l4-cd}

\section{The absence of expansion code}
The token reader skips whitespace, recognizes operator characters and collects
ordinary word spans. It has no state for single or double quotes, and no stage
that expands variables or enumerates wildcard matches. Searching for an asterisk
finds C pointer syntax, not an implementation of globbing; searching for a dollar
sign finds the prompt. Those search hits require interpretation in their surrounding
code. This source reading agrees with the literal stars, dollar-names and quote
characters measured by the argument program.

\utmconclusions
The completed build and boot establish a working xv6 environment on openSUSE.
The sleep utility reuses the already-correct pause interface from Lab 2, and the
argument utility isolates the shell's interpretation from application behaviour.
Both shells rely on familiar process and descriptor mechanisms, but Bash performs
substantial additional work before a program starts. The source analysis explains
both the missing expansion features and the narrower treatment of directory changes.
All prescribed command experiments were executed, and QEMU was shut down through
its escape command so the filesystem image is released for subsequent use.
\end{document}
